
\makeatletter
\def\input@path{{./}{../}{../../}{../../../}{../../../../}{preamble/}{../preamble/}{../../preamble/}{../../../preamble/}}
\makeatother

\input{preamble.tex}
\input{graphs/highlights.tex}
% \GTenHighlightsfalse % Uncomment to disable instructional highlights.
\input{graphs/specific_highlights.tex}
\COneTonePracticePalette

\input{graphs/uniform_linear_probability_distributions.tex}
\input{graphs/probability_distribution_graphs.tex}
\input{graphs/normal_distribution_graphs.tex}
\input{graphs/G11_C5_graphs.tex}
\usepackage{longtable}


\setlength{\DistributionGraphWidth}{0.6\textwidth}
\setlength{\DistributionGraphHeight}{6cm}

\setlength{\DistributionGraphHeightCompact}{3.5cm}

\title{Grade 11 -- Term 1 \\ C5 Practice Activity 5 with Solutions}
\author{}
\date{}

\pagestyle{empty}

% ============================================================
% GRAPH WRAPPER
% ============================================================

\newcommand{\SolvedGraph}[2]{%
	\par\smallskip
	\Needspace{4.6cm}%
	\begin{center}
		\textbf{\small #1}\par\vspace{-1mm}
		\begin{minipage}{\textwidth}
			#2
		\end{minipage}%
	\end{center}
	\vspace{-2mm}%
}

% Each solution stage is numbered as Problem.Stage.
\newcommand{\SolutionStage}[2]{%
	\Needspace{6\baselineskip}%
	\subsubsection{#2}%
}

\begin{document}

	\maketitle
	\setcounter{tocdepth}{1}
	\setcounter{secnumdepth}{2}
	\phantomsection\label{toc}
	\vspace*{-20mm}
	\tableofcontents
	
	\section*{Evaluated criterion}
	\begin{itemize}
		\item [\textbf{C5.}] Calculates and interprets the mean and variance of continuous distributions using standard formulas and functions.
	\end{itemize}
	
	\vspace{-4mm}
	
	\section{Introduction}\label{sec:c4-introduction}
\Needspace{8\baselineskip}
\subsection*{Mathematical highlight guide}
\phantomsection\label{tone:C5:legend}
\PracticeHighlightLegend{G11-T1-C5}

	\SectionProblemsTableOfContents{
		\SectionProblemLink{sec:c4-intro-purpose}{Purpose and learning goals} \\
		\SectionProblemLink{sec:c4-intro-distributions}{Continuous probability distributions} \\
		\SectionProblemLink{sec:c4-intro-method}{Mean, variance, and standard deviation} \\
		\SectionProblemLink{sec:c4-intro-formulas}{Formula reference}
	}\vspace{8mm}
	
	Criterion C4 focuses on calculating and interpreting measures that describe the
	center and spread of a continuous probability distribution.  The main measures
	used are the mean, or \CfiveToneExpected{expected} \CfiveToneExpected{value}, the variance, and the standard deviation.
	
	These measures summarize different features of a probability distribution.  The
	mean identifies its center, while the variance and standard deviation describe
	how widely values are distributed around that center.  The appropriate formulas
	depend on the probability distribution and its parameters.
	
\subsection{Purpose and learning goals}\label{sec:c4-intro-purpose}
The purpose of this activity is to calculate and interpret the mean, variance,
and standard deviation of continuous probability distributions using standard
formulas.

\begin{center}
	\renewcommand{\arraystretch}{1.2}
	\begin{tabular}{lll}
		\toprule
		Measure & Common symbols & Meaning \\
		\midrule
		Mean
		& $\CfiveToneExpected{E}(\ConeToneVariable{X})$, $\CsevenToneSampleMean{\bar{x}}$, $\CfiveTonePopulationMean{\mu}$
		& Expected or average value \\
		
		Variance
		& $V\ConeToneVariable{a}\ConeToneVariable{r}(\ConeToneVariable{X})$, $\CsevenToneSampleVariance{s^2}$, $\CfiveTonePopulationVariance{\sigma^2}$
		& Average squared spread from the mean \\
		
		Standard deviation
		& $\ConeToneVariable{S}\ConeToneVariable{D}(\ConeToneVariable{X})$, $\CsevenToneSampleSD{s}$, $\CfiveTonePopulationSD{\sigma}$
		& Typical spread from the mean \\
		\bottomrule
	\end{tabular}
\end{center}

For each situation:
\begin{enumerate}
	\item Identify the probability distribution and its parameters.
	\item Calculate the mean, or \CfiveToneExpected{expected} \CfiveToneExpected{value}, $\CfiveToneExpected{E}(\ConeToneVariable{X})$.
	\item Calculate the variance, $V\ConeToneVariable{a}\ConeToneVariable{r}(\ConeToneVariable{X})$.
	\item Calculate the standard deviation,
	$\CfiveTonePopulationSD{\sigma_X}=\sqrt{V\ConeToneVariable{a}\ConeToneVariable{r}(\ConeToneVariable{X})}$.
	\item Interpret the results in the context of the variable.
\end{enumerate}

	\secInicio[sec:c4-intro-purpose][sec:c4-introduction]
	
	\subsection{Continuous probability distributions}\label{sec:c4-intro-distributions}
	
	The distributions used in this criterion describe continuous variables.  Their
	probability density functions assign probability to intervals through area
	under the curve.  Different distributions have different shapes and parameters,
	so their means and variances are calculated using different formulas.
	
	The distributions used in this activity are:
	\begin{itemize}
		\item Uniform
		\item Triangular
		\item Linear increasing
		\item Linear decreasing
		\item Trapezoidal
		\item Exponential
		\item Normal
	\end{itemize}
	
	The notation identifies the distribution and its parameters.  For example,
	\[
	\ConeToneVariable{X}\sim \ConeToneVariable{U}(\ConeToneVariable{a},\ConeToneVariable{b})
	\]
	indicates a uniform distribution between $\ConeToneVariable{a}$ and $\ConeToneVariable{b}$, while
	\[
	\ConeToneVariable{X}\sim N(\CfiveTonePopulationMean{\mu},\CfiveTonePopulationVariance{\sigma^2})
	\]
	indicates a normal distribution with mean $\CfiveTonePopulationMean{\mu}$ and variance $\CfiveTonePopulationVariance{\sigma^2}$.
	\secInicio[sec:c4-intro-distributions][sec:c4-introduction]
	
	\subsection{Mean, variance, and standard deviation}\label{sec:c4-intro-method}
	
	The mean, written as $\CfiveToneExpected{E}(\ConeToneVariable{X})$, describes the center of a probability distribution.
	For a continuous \ConeToneVariable{random} \ConeToneVariable{variable}, it represents the long-run average value
	expected from repeated observations under the model.
	
	Variance measures the spread of the distribution around its mean:
	\[
	V\ConeToneVariable{a}\ConeToneVariable{r}(\ConeToneVariable{X})=\CfiveToneExpected{E}\left[(\ConeToneVariable{X}-\CfiveToneExpected{E}(\ConeToneVariable{X}))^2\right].
	\]
	The variance is expressed in squared units.  Its square root is the standard
	deviation:
	\[
	\CfiveTonePopulationSD{\sigma_X}=\sqrt{V\ConeToneVariable{a}\ConeToneVariable{r}(\ConeToneVariable{X})}.
	\]
	Unlike variance, the standard deviation has the same units as the original
	variable, which makes it easier to interpret in context.
	
	For every problem, first identify the distribution and its parameters.  Then
	select the corresponding formulas, substitute the parameters, calculate the
	results, and interpret their meaning using the units of the variable.
	\secInicio[sec:c4-intro-method][sec:c4-introduction]
	
	\subsection{Formula reference}\label{sec:c4-intro-formulas}
	
	The following table summarizes the formulas used in this criterion.
	
	{
		\renewcommand{\arraystretch}{1.9}
		
		\begin{longtable}{c|c|c|c}
			$\text{Distribution}$
			&
			$\text{Notation}$
			&
			$\CfiveToneExpected{E}(\ConeToneVariable{X})$
			&
			$V\ConeToneVariable{a}\ConeToneVariable{r}(\ConeToneVariable{X})$
			\\[1mm]
			\hline
			\endfirsthead
			
			$\text{Distribution}$
			&
			$\text{Notation}$
			&
			$\CfiveToneExpected{E}(\ConeToneVariable{X})$
			&
			$V\ConeToneVariable{a}\ConeToneVariable{r}(\ConeToneVariable{X})$
			\\[1mm]
			\hline
			\endhead
			
			\hline
			\endfoot
			
			$\text{Uniform}$
			&
			$\ConeToneVariable{X}\sim \ConeToneVariable{U}(\ConeToneVariable{a},\ConeToneVariable{b})$
			&
			$\dfrac{\ConeToneVariable{a}+\ConeToneVariable{b}}{2}$
			&
			$\dfrac{(\ConeToneVariable{b}-\ConeToneVariable{a})^2}{12}$
			\\[2mm]
			
			&
			\multicolumn{3}{c}{
				$\vcenter{\hbox{\CompactUniformGraph{0}{6}}}$
			}
			\\[5mm]
			\hline
			
			$\text{Triangular}$
			&
			$\ConeToneVariable{X}\sim \operatorname{Tri}(\ConeToneVariable{a},\ConeToneVariable{b},\CtwoToneMode{c})$
			&
			$\dfrac{\ConeToneVariable{a}+\ConeToneVariable{b}+\CtwoToneMode{c}}{3}$
			&
			$\dfrac{a^2+b^2+c^2-\ConeToneVariable{a}\ConeToneVariable{b}-\ConeToneVariable{a}\CtwoToneMode{c}-\ConeToneVariable{b}\CtwoToneMode{c}}{18}$
			\\[2mm]
			
			&
			\multicolumn{3}{c}{
				$\vcenter{\hbox{\CompactTriangularGraph{0}{3}{6}}}$
			}
			\\[5mm]
			\hline
			
			$\text{Linear increasing}$
			&
			$\ConeToneVariable{X}\sim \operatorname{Lin}_{\mathrm{inc}}(\ConeToneVariable{a},\ConeToneVariable{b})$
			&
			$\dfrac{\ConeToneVariable{a}+\ConeToneVariable{2}\ConeToneVariable{b}}{3}$
			&
			$\dfrac{(\ConeToneVariable{b}-\ConeToneVariable{a})^2}{18}$
			\\[2mm]
			
			&
			\multicolumn{3}{c}{
				$\vcenter{\hbox{\CompactLinearIncreasingGraph{0}{6}}}$
			}
			\\[5mm]
			\hline
			
			$\text{Linear decreasing}$
			&
			$\ConeToneVariable{X}\sim \operatorname{Lin}_{\mathrm{dec}}(\ConeToneVariable{a},\ConeToneVariable{b})$
			&
			$\dfrac{\ConeToneVariable{2}\ConeToneVariable{a}+\ConeToneVariable{b}}{3}$
			&
			$\dfrac{(\ConeToneVariable{b}-\ConeToneVariable{a})^2}{18}$
			\\[2mm]
			
			&
			\multicolumn{3}{c}{
				$\vcenter{\hbox{\CompactLinearDecreasingGraph{0}{6}}}$
			}
			\\[5mm]
			\hline
			
			\newpage
			
			$\text{Trapezoidal}$
			&
			$\ConeToneVariable{X}\sim \operatorname{Trap}(\ConeToneVariable{a},\ConeToneVariable{b},\CtwoToneDensity{h_1},\CtwoToneDensity{h_2})$
			&
			$\dfrac{\ConeToneVariable{b}-\ConeToneVariable{a}}{3}\,
			\dfrac{\CtwoToneDensity{h_1^2}+\CtwoToneDensity{h_1}\CtwoToneDensity{h_2}+\CtwoToneDensity{h_2^2}}{\CtwoToneDensity{h_1}+\CtwoToneDensity{h_2}}
			+\dfrac{\ConeToneVariable{a}+\ConeToneVariable{b}}{2}$
			&
			$\dfrac{(\ConeToneVariable{b}-\ConeToneVariable{a})^2}{18}\,
			\dfrac{\CtwoToneDensity{h_1^2}+\ConeToneVariable{2}\CtwoToneDensity{h_1}\CtwoToneDensity{h_2}+\CtwoToneDensity{h_2^2}}{(\CtwoToneDensity{h_1}+\CtwoToneDensity{h_2})^2}$
			\\[2mm]
			
			&
			\multicolumn{3}{c}{
				$\vcenter{\hbox{
						\CompactTrapezoidalGraph{0}{6}{0.1111111}{0.3333333}
				}}$
			}
			\\[5mm]
			\hline
			
			$\text{Exponential}$
			&
			$\ConeToneVariable{X}\sim \operatorname{Exp}(\CfourToneRate{\lambda})$
			&
			$\dfrac{\ConeToneVariable{1}}{\CfourToneRate{\lambda}}$
			&
			$\dfrac{\ConeToneVariable{1}}{\CfourToneRate{\lambda^2}}$
			\\[2mm]
			
			&
			\multicolumn{3}{c}{
				$\vcenter{\hbox{\CompactExponentialGraph{1}}}$
			}
			\\[5mm]
			\hline
			
			$\text{Normal}$
			&
			$\ConeToneVariable{X}\sim N(\CfiveTonePopulationMean{\mu},\CfiveTonePopulationVariance{\sigma^2})$
			&
			$\CfiveTonePopulationMean{\mu}$
			&
			$\CfiveTonePopulationVariance{\sigma^2}$
			\\[2mm]
			
			&
			\multicolumn{3}{c}{
				$\vcenter{\hbox{\CompactNormalGraph{0}{1}}}$
			}
			
		\end{longtable}
	}
	\secInicio[sec:c4-intro-formulas][sec:c4-introduction]
	
	The standard deviation is obtained after calculating the variance:
	\[
	\CfiveTonePopulationSD{\sigma_X}=\sqrt{V\ConeToneVariable{a}\ConeToneVariable{r}(\ConeToneVariable{X})}.
	\]
	When interpreting a result, report the mean and standard deviation using the
	units of the original variable, while variance uses squared units.
	
	The formulas above should be applied directly once the distribution and its
	parameters have been identified.  The second table provides a compact
	reference without the graphs.
	
	{
		\renewcommand{\arraystretch}{1.9}
		
		\begin{longtable}{c|c|c|c}
			$\text{Distribution}$
			&
			$\text{Notation}$
			&
			$\CfiveToneExpected{E}(\ConeToneVariable{X})$
			&
			$V\ConeToneVariable{a}\ConeToneVariable{r}(\ConeToneVariable{X})$
			\\[1mm]
			\hline
			\endfirsthead
			
			$\text{Distribution}$
			&
			$\text{Notation}$
			&
			$\CfiveToneExpected{E}(\ConeToneVariable{X})$
			&
			$V\ConeToneVariable{a}\ConeToneVariable{r}(\ConeToneVariable{X})$
			\\[1mm]
			\hline
			\endhead
			
			\hline
			\endfoot
			
			$\text{Uniform}$
			&
			$\ConeToneVariable{X}\sim \ConeToneVariable{U}(\ConeToneVariable{a},\ConeToneVariable{b})$
			&
			$\dfrac{\ConeToneVariable{a}+\ConeToneVariable{b}}{2}$
			&
			$\dfrac{(\ConeToneVariable{b}-\ConeToneVariable{a})^2}{12}$
			\\[3mm]
			\hline
			
			$\text{Triangular}$
			&
			$\ConeToneVariable{X}\sim \operatorname{Tri}(\ConeToneVariable{a},\ConeToneVariable{b},\CtwoToneMode{c})$
			&
			$\dfrac{\ConeToneVariable{a}+\ConeToneVariable{b}+\CtwoToneMode{c}}{3}$
			&
			$\dfrac{a^2+b^2+c^2-\ConeToneVariable{a}\ConeToneVariable{b}-\ConeToneVariable{a}\CtwoToneMode{c}-\ConeToneVariable{b}\CtwoToneMode{c}}{18}$
			\\[3mm]
			\hline
			
			$\text{Linear increasing}$
			&
			$\ConeToneVariable{X}\sim \operatorname{Lin}_{\mathrm{inc}}(\ConeToneVariable{a},\ConeToneVariable{b})$
			&
			$\dfrac{\ConeToneVariable{a}+\ConeToneVariable{2}\ConeToneVariable{b}}{3}$
			&
			$\dfrac{(\ConeToneVariable{b}-\ConeToneVariable{a})^2}{18}$
			\\[3mm]
			\hline
			
			$\text{Linear decreasing}$
			&
			$\ConeToneVariable{X}\sim \operatorname{Lin}_{\mathrm{dec}}(\ConeToneVariable{a},\ConeToneVariable{b})$
			&
			$\dfrac{\ConeToneVariable{2}\ConeToneVariable{a}+\ConeToneVariable{b}}{3}$
			&
			$\dfrac{(\ConeToneVariable{b}-\ConeToneVariable{a})^2}{18}$
			\\[3mm]
			\hline
			
			$\text{Trapezoidal}$
			&
			$\ConeToneVariable{X}\sim \operatorname{Trap}(\ConeToneVariable{a},\ConeToneVariable{b},\CtwoToneDensity{h_1},\CtwoToneDensity{h_2})$
			&
			$\dfrac{\ConeToneVariable{b}-\ConeToneVariable{a}}{3}\,
			\dfrac{\CtwoToneDensity{h_1^2}+\CtwoToneDensity{h_1}\CtwoToneDensity{h_2}+\CtwoToneDensity{h_2^2}}{\CtwoToneDensity{h_1}+\CtwoToneDensity{h_2}}
			+\dfrac{\ConeToneVariable{a}+\ConeToneVariable{b}}{2}$
			&
			$\dfrac{(\ConeToneVariable{b}-\ConeToneVariable{a})^2}{18}\,
			\dfrac{\CtwoToneDensity{h_1^2}+\ConeToneVariable{2}\CtwoToneDensity{h_1}\CtwoToneDensity{h_2}+\CtwoToneDensity{h_2^2}}{(\CtwoToneDensity{h_1}+\CtwoToneDensity{h_2})^2}$
			\\[3mm]
			\hline
			
			$\text{Exponential}$
			&
			$\ConeToneVariable{X}\sim \operatorname{Exp}(\CfourToneRate{\lambda})$
			&
			$\dfrac{\ConeToneVariable{1}}{\CfourToneRate{\lambda}}$
			&
			$\dfrac{\ConeToneVariable{1}}{\CfourToneRate{\lambda^2}}$
			\\[3mm]
			\hline
			
			$\text{Normal}$
			&
			$\ConeToneVariable{X}\sim N(\CfiveTonePopulationMean{\mu},\CfiveTonePopulationVariance{\sigma^2})$
			&
			$\CfiveTonePopulationMean{\mu}$
			&
			$\CfiveTonePopulationVariance{\sigma^2}$
			
		\end{longtable}
	}
	\secInicio[sec:c4-intro-formulas][sec:c4-introduction]
	
	
	\newpage
	
\section{Introductory Problems}\label{sec:c4-problems}

\SectionProblemsTableOfContents{
	\SectionProblemLink{sec:c4-problem-uniform}{Uniform distribution} \\
	\SectionProblemLink{sec:c4-problem-linear-increasing}{Linear increasing distribution} \\
	\SectionProblemLink{sec:c4-problem-linear-decreasing}{Linear decreasing distribution} \\
	\SectionProblemLink{sec:c4-problem-triangular}{Triangular distribution} \\
	\SectionProblemLink{sec:c4-problem-exponential}{Exponential distribution} \\
	\SectionProblemLink{sec:c4-problem-normal}{Normal distribution}
}

\subsection{Uniform distribution}\label{sec:c4-problem-uniform}

\SolutionStage{1}{Distribution and parameters}

A small business invests in a short-term opportunity. Based on similar
investments, the money earned is expected to be between \$2,000 and
\$12,000, with all values in this range considered equally likely.

Let $\ConeToneVariable{X}$ represent the money earned, measured in thousands of dollars.
Therefore,
\[
\ConeToneVariable{X}\sim \ConeToneVariable{U}(\ConeToneVariable{a},\ConeToneVariable{b})
\qquad \longrightarrow \qquad
\ConeToneVariable{X}\sim \ConeToneVariable{U}(2,12)
\qquad \longrightarrow \qquad
\ConeToneVariable{a}=\ConeToneVariable{2},\qquad \ConeToneVariable{b}=\ConeToneVariable{12}.
\]

\SolvedGraph{Uniform density}{
	\UniformGraph{2}{12}
}

\SolutionStage{2}{Mean}

For a uniform distribution,
\[
\CfiveToneExpected{E}(\ConeToneVariable{X})=\frac{\ConeToneVariable{a}+\ConeToneVariable{b}}{2}
\qquad \longrightarrow \qquad
\CfiveToneExpected{\CfiveToneExpected{E}}(\ConeToneVariable{\ConeToneVariable{X}})=\frac{\ConeToneVariable{\ConeToneVariable{2}}+\ConeToneVariable{\ConeToneVariable{12}}}{2}
=7
\qquad \longrightarrow \qquad
\boxed{\CfiveToneExpected{E}(\ConeToneVariable{X})=\CfiveToneExpected{7}\text{ thousand dollars}}
\]

\SolutionStage{3}{Variance}

For a uniform distribution,
\[
V\ConeToneVariable{a}\ConeToneVariable{r}(\ConeToneVariable{X})=\frac{(\ConeToneVariable{b}-\ConeToneVariable{a})^2}{12}
\qquad \longrightarrow \qquad
V\ConeToneVariable{a}\ConeToneVariable{r}(\ConeToneVariable{X})
=\frac{(\ConeToneVariable{12}\ConeToneVariable{-2})^2}{12}
=\frac{\ConeToneVariable{100}}{12}
=\frac{\ConeToneVariable{25}}{3}
\approx\ConeToneVariable{8.33}
\qquad \longrightarrow \qquad
\boxed{V\ConeToneVariable{a}\ConeToneVariable{r}(\ConeToneVariable{X})\approx\ConeToneVariable{8.33}\text{ thousand dollars}^2}
\]

\SolutionStage{4}{Standard Deviation}

The standard deviation is
\[
\CfiveTonePopulationSD{\sigma_X}=\sqrt{\frac{\ConeToneVariable{25}}{3}}
\approx\CfiveTonePopulationSD{2.89}
\qquad \longrightarrow \qquad
\boxed{\CfiveTonePopulationSD{\sigma_X}\approx\CfiveTonePopulationSD{2.89}\text{ thousand dollars}}
\]

\SolutionStage{5}{Probability within one standard deviation}

We want the probability that the money earned is within one standard
deviation of the mean:
\[
\ConeToneProbability{P}(\CfiveTonePopulationMean{\mu}-\CfiveTonePopulationSD{\sigma_X}<\ConeToneVariable{X}<\CfiveTonePopulationMean{\mu}+\CfiveTonePopulationSD{\sigma_X})
\]

Using $\CfiveTonePopulationMean{\mu}=\CfiveTonePopulationMean{7}$ and $\CfiveTonePopulationSD{\sigma_X}\approx\CfiveTonePopulationSD{2.89}$,
\[
\ConeToneProbability{\ConeToneProbability{P}}(\CfiveTonePopulationMean{7}-\CfiveTonePopulationSD{2.89}<\ConeToneVariable{\ConeToneVariable{X}}<\CfiveTonePopulationMean{7}+\CfiveTonePopulationSD{2.89})
=
\ConeToneProbability{P}(\ConeToneProbability{4.11}<\ConeToneVariable{X}<9.89)
\]

For a uniform distribution,
\[
\ConeToneProbability{P}(\CtwoToneMode{c}<\ConeToneVariable{X}<\ConeToneVariable{d})=\frac{\ConeToneVariable{d}-\CtwoToneMode{c}}{\ConeToneVariable{b}-\ConeToneVariable{a}}
\]

Therefore,
\[
\ConeToneProbability{P}(4.11<\ConeToneVariable{X}<9.89)
=
\frac{\ConeToneVariable{9.89}\ConeToneVariable{-4.11}}{12-2}
=
\frac{\ConeToneVariable{5.78}}{10}
\approx\ConeToneProbability{0.578}
\]

\[
\boxed{\ConeToneProbability{P}(\CfiveTonePopulationMean{\mu}-\CfiveTonePopulationSD{\sigma_X}<\ConeToneVariable{X}<\CfiveTonePopulationMean{\mu}+\CfiveTonePopulationSD{\sigma_X})\approx\ConeToneProbability{0.578}}
\]

Thus, approximately $57.8\%$ of the possible investment earnings
are within one standard deviation of the mean.

\SolutionStage{6}{Probability within one standard deviation for any uniform distribution}

For a uniform distribution $\ConeToneVariable{X}\sim \ConeToneVariable{U}(\ConeToneVariable{a},\ConeToneVariable{b})$, the mean and standard
deviation are
\[
\CfiveTonePopulationMean{\mu}=\frac{\ConeToneVariable{a}+\ConeToneVariable{b}}{2}
\qquad\text{and}\qquad
\CfiveTonePopulationSD{\sigma_X}=\frac{\ConeToneVariable{b}-\ConeToneVariable{a}}{\sqrt{12}}.
\]

The interval within one standard deviation of the mean is
\[
\CfiveTonePopulationMean{\mu}-\CfiveTonePopulationSD{\sigma_X}<\ConeToneVariable{X}<\CfiveTonePopulationMean{\mu}+\CfiveTonePopulationSD{\sigma_X}.
\]

Its total width is
\[
(\CfiveTonePopulationMean{\mu}+\CfiveTonePopulationSD{\sigma_X})-(\CfiveTonePopulationMean{\mu}-\CfiveTonePopulationSD{\sigma_X})
=\CfiveTonePopulationMean{2}\CfiveTonePopulationSD{\sigma_X}
=\frac{\ConeToneVariable{2}(\ConeToneVariable{b}-\ConeToneVariable{a})}{\sqrt{12}}.
\]

Since the probability of an interval in a uniform distribution is its
width divided by the total width $\ConeToneVariable{b}-\ConeToneVariable{a}$,
\[
\ConeToneProbability{P}(\CfiveTonePopulationMean{\mu}-\CfiveTonePopulationSD{\sigma_X}<\ConeToneVariable{X}<\CfiveTonePopulationMean{\mu}+\CfiveTonePopulationSD{\sigma_X})
=
\frac{\ConeToneVariable{2}\CfiveTonePopulationSD{\sigma_X}}{\ConeToneVariable{b}-\ConeToneVariable{a}}
=
\frac{\ConeToneVariable{2}}{\sqrt{12}}
=
\frac{\ConeToneVariable{1}}{\sqrt{3}}
\approx\ConeToneProbability{0.577}.
\]

Therefore, for every uniform distribution,
\[
\boxed{
	\ConeToneProbability{P}(\CfiveTonePopulationMean{\mu}-\CfiveTonePopulationSD{\sigma_X}<\ConeToneVariable{X}<\CfiveTonePopulationMean{\mu}+\CfiveTonePopulationSD{\sigma_X})
	=\frac{\ConeToneVariable{1}}{\sqrt{3}}
	\approx\ConeToneProbability{57.7\%}
}
\]

Thus, the probability within one standard deviation of the mean is
the same for every continuous uniform distribution.

\subsection{Linear increasing distribution}\label{sec:c4-problem-linear-increasing}

\SolutionStage{1}{Distribution and parameters}

A bakery prepares custom cakes for weekend orders. Based on previous
weekends, the preparation time for a randomly selected cake ranges from
80 to 140 minutes. Longer preparation times are more likely because
larger and more detailed cakes require more work, so the probability
density increases linearly across this range.

Let $\ConeToneVariable{T}$ represent the preparation time, in minutes, for one cake. Therefore,
\[
\ConeToneVariable{T}\sim\operatorname{Lin}_{\mathrm{inc}}(\ConeToneVariable{a},\ConeToneVariable{b})
\qquad \longrightarrow \qquad
\ConeToneVariable{T}\sim\operatorname{Lin}_{\mathrm{inc}}(80,140)
\qquad \longrightarrow \qquad
\ConeToneVariable{a}=\ConeToneVariable{80}\qquad \ConeToneVariable{b}=\ConeToneVariable{140}
\]

\SolvedGraph{Linear increasing density}{
	\LinearIncreasingGraph{80}{140}
}

\SolutionStage{2}{Mean}

For a linear increasing distribution,
\[
\CfiveToneExpected{E}(\ConeToneVariable{T})=\frac{\ConeToneVariable{a}+\ConeToneVariable{2}\ConeToneVariable{b}}{3}
\qquad \longrightarrow \qquad
\CfiveToneExpected{E}(\ConeToneVariable{T})
=\frac{\ConeToneVariable{80}+\ConeToneVariable{2}(\ConeToneVariable{140})}{3}
=\frac{\ConeToneVariable{360}}{3}
=\CfiveToneExpected{120}
\qquad \longrightarrow \qquad
\boxed{\CfiveToneExpected{E}(\ConeToneVariable{T})=\CfiveToneExpected{120}\text{ minutes}}
\]

\SolutionStage{3}{Variance}

For a linear increasing distribution,
\[
V\ConeToneVariable{a}\ConeToneVariable{r}(\ConeToneVariable{T})=\frac{(\ConeToneVariable{b}-\ConeToneVariable{a})^2}{18}
\qquad \longrightarrow \qquad
V\ConeToneVariable{a}\ConeToneVariable{r}(\ConeToneVariable{T})
=\frac{(\ConeToneVariable{140}\ConeToneVariable{-80})^2}{18}
=\frac{\ConeToneVariable{3600}}{18}
=\ConeToneVariable{200}
\qquad \longrightarrow \qquad
\boxed{V\ConeToneVariable{a}\ConeToneVariable{r}(\ConeToneVariable{T})=\ConeToneVariable{200}\text{ minutes}^2}
\]

\SolutionStage{4}{Interpretation}

The standard deviation is
\[
\CfiveTonePopulationSD{\sigma_T}=\sqrt{\CfiveTonePopulationVariance{200}}
\approx\CfiveTonePopulationSD{14.14}
\qquad \longrightarrow \qquad
\boxed{\CfiveTonePopulationSD{\sigma_T}\approx\CfiveTonePopulationSD{14.14}\text{ minutes}}
\]

\SolutionStage{5}{Probability within one standard deviation}

We want the probability that the preparation time is within one standard
deviation of the mean:
\[
\ConeToneProbability{P}(\CfiveTonePopulationMean{\mu}-\CfiveTonePopulationSD{\sigma_T}<\ConeToneVariable{T}<\CfiveTonePopulationMean{\mu}+\CfiveTonePopulationSD{\sigma_T})
\]

Using $\CfiveTonePopulationMean{\mu}=\CfiveTonePopulationMean{120}$ and $\CfiveTonePopulationSD{\sigma_T}\approx\CfiveTonePopulationSD{14.14}$,
\[
\ConeToneProbability{P}(120-14.14<\ConeToneVariable{T}<120+14.14)
=
\ConeToneProbability{P}(\ConeToneProbability{105.86}<\ConeToneVariable{T}<134.14)
\]

For a linear increasing distribution, the probability density is
\[
\CtwoToneDensity{f}(\ConeToneVariable{t})=\frac{\ConeToneVariable{2}(\ConeToneVariable{t}-\ConeToneVariable{a})}{(\ConeToneVariable{b}-\ConeToneVariable{a})^2}.
\qquad \longrightarrow \qquad
\begin{aligned}
	\ConeToneProbability{P}(105.86<\ConeToneVariable{T}<134.14)
	&=\int_{105.86}^{134.14}
	\frac{\ConeToneVariable{2}(\ConeToneVariable{t}\ConeToneVariable{-80})}{(140-80)^2}\,\ConeToneVariable{d}\ConeToneVariable{t}\\[\ConeToneProbability{4}\ConeToneProbability{p}\ConeToneVariable{t}]
	&=\left[
	\frac{(\ConeToneVariable{t}\ConeToneVariable{-80})^2}{3600}
	\right]_{105.86}^{134.14}\\[\ConeToneProbability{4}\ConeToneProbability{p}\ConeToneVariable{t}]
	&\approx \ConeToneProbability{0.6285}
\end{aligned}
\]

\[
\boxed{\ConeToneProbability{P}(\CfiveTonePopulationMean{\mu}-\CfiveTonePopulationSD{\sigma_T}<\ConeToneVariable{T}<\CfiveTonePopulationMean{\mu}+\CfiveTonePopulationSD{\sigma_T})\approx\ConeToneProbability{0.6285}}
\]

Thus, approximately $62.84\%$ of cake preparation times are within one
standard deviation of the mean.


\subsection{Linear decreasing distribution}\label{sec:c4-problem-linear-decreasing}

\SolutionStage{1}{Distribution and parameters}

A bakery delivers custom cake orders to customers during the weekend.
Based on previous deliveries, the delivery time for a randomly selected
order ranges from 20 to 80 minutes. Shorter delivery times are more likely
because most orders are delivered within the local area, so the probability
density decreases linearly as the delivery time increases.

Let $\ConeToneVariable{T}$ represent the delivery time, in minutes, for one order. Therefore,
\[
\ConeToneVariable{T}\sim\operatorname{Lin}_{\mathrm{dec}}(\ConeToneVariable{a},\ConeToneVariable{b})
\qquad \longrightarrow \qquad
\ConeToneVariable{T}\sim\operatorname{Lin}_{\mathrm{dec}}(20,80)
\qquad \longrightarrow \qquad
\ConeToneVariable{a}=\ConeToneVariable{20},\qquad \ConeToneVariable{b}=\ConeToneVariable{80}
\]

\SolvedGraph{Linear decreasing density}{
	\LinearDecreasingGraph{20}{80}
}

\SolutionStage{2}{Mean}

For a linear decreasing distribution,
\[
\CfiveToneExpected{E}(\ConeToneVariable{T})=\frac{\ConeToneVariable{2}\ConeToneVariable{a}+\ConeToneVariable{b}}{3}
\qquad \longrightarrow \qquad
\CfiveToneExpected{E}(\ConeToneVariable{T})
=\frac{\ConeToneVariable{2}(\ConeToneVariable{20})+\ConeToneVariable{80}}{3}
=\frac{\ConeToneVariable{120}}{3}
=\CfiveToneExpected{40}
\qquad \longrightarrow \qquad
\boxed{\CfiveToneExpected{E}(\ConeToneVariable{T})=\CfiveToneExpected{40}\text{ minutes}}
\]

\SolutionStage{3}{Variance}

For a linear decreasing distribution,
\[
V\ConeToneVariable{a}\ConeToneVariable{r}(\ConeToneVariable{T})=\frac{(\ConeToneVariable{b}-\ConeToneVariable{a})^2}{18}
\qquad \longrightarrow \qquad
V\ConeToneVariable{a}\ConeToneVariable{r}(\ConeToneVariable{T})
=\frac{(\ConeToneVariable{80}\ConeToneVariable{-20})^2}{18}
=\frac{\ConeToneVariable{3600}}{18}
=\ConeToneVariable{200}
\qquad \longrightarrow \qquad
\boxed{V\ConeToneVariable{a}\ConeToneVariable{r}(\ConeToneVariable{T})=\ConeToneVariable{200}\text{ minutes}^2}
\]

\SolutionStage{4}{Interpretation}

The standard deviation is
\[
\CfiveTonePopulationSD{\sigma_T}=\sqrt{\CfiveTonePopulationVariance{200}}
\approx\CfiveTonePopulationSD{14.14}
\qquad \longrightarrow \qquad
\boxed{\CfiveTonePopulationSD{\sigma_T}\approx\CfiveTonePopulationSD{14.14}\text{ minutes}}
\]

\SolutionStage{5}{Probability within one standard deviation}

We want the probability that the delivery time is within one standard
deviation of the mean:
\[
\ConeToneProbability{P}(\CfiveTonePopulationMean{\mu}-\CfiveTonePopulationSD{\sigma_T}<\ConeToneVariable{T}<\CfiveTonePopulationMean{\mu}+\CfiveTonePopulationSD{\sigma_T})
\]

Using $\CfiveTonePopulationMean{\mu}=\CfiveTonePopulationMean{40}$ and $\CfiveTonePopulationSD{\sigma_T}\approx\CfiveTonePopulationSD{14.14}$,
\[
\ConeToneProbability{P}(40-14.14<\ConeToneVariable{T}<40+14.14)
=
\ConeToneProbability{P}(\ConeToneProbability{25.86}<\ConeToneVariable{T}<54.14)
\]

For a linear decreasing distribution, the probability density is
\[
\CtwoToneDensity{f}(\ConeToneVariable{t})=\frac{\ConeToneVariable{2}(\ConeToneVariable{b}-\ConeToneVariable{t})}{(\ConeToneVariable{b}-\ConeToneVariable{a})^2}.
\qquad \longrightarrow \qquad 
\begin{aligned}
	\ConeToneProbability{P}(25.86<\ConeToneVariable{T}<54.14)
	&=\int_{25.86}^{54.14}
	\frac{\ConeToneVariable{2}(\ConeToneVariable{80}-\ConeToneVariable{t})}{(80-20)^2}\,\ConeToneVariable{d}\ConeToneVariable{t}\\[\ConeToneProbability{4}\ConeToneProbability{p}\ConeToneVariable{t}]
	&=\left[
	\frac{\ConeToneVariable{2}(\ConeToneVariable{80})\ConeToneVariable{t}-t^2}{3600}
	\right]_{25.86}^{54.14}\\[\ConeToneProbability{4}\ConeToneProbability{p}\ConeToneVariable{t}]
	&\approx \ConeToneProbability{0.6285}
\end{aligned}
\]
\[
\boxed{\ConeToneProbability{P}(\CfiveTonePopulationMean{\mu}-\CfiveTonePopulationSD{\sigma_T}<\ConeToneVariable{T}<\CfiveTonePopulationMean{\mu}+\CfiveTonePopulationSD{\sigma_T})\approx\ConeToneProbability{0.6285}}
\]

Thus, approximately $62.84\%$ of delivery times are within one
standard deviation of the mean.

\subsubsection{Probability within one standard deviation}

The probability of obtaining a value within one standard deviation of the
mean is the same for every linear increasing or decreasing distribution,
regardless of the values of $\ConeToneVariable{a}$ and $\ConeToneVariable{b}$.

For a linear increasing distribution,
\[
\CtwoToneDensity{f}(\ConeToneVariable{x})=\frac{\ConeToneVariable{2}(\ConeToneVariable{x}-\ConeToneVariable{a})}{(\ConeToneVariable{b}-\ConeToneVariable{a})^2},
\qquad
\CfiveTonePopulationMean{\mu}=\frac{\ConeToneVariable{a}+\ConeToneVariable{2}\ConeToneVariable{b}}{3},
\qquad
\CfiveTonePopulationSD{\sigma}=\frac{\ConeToneVariable{b}-\ConeToneVariable{a}}{\sqrt{18}},
\qquad \ConeToneVariable{a}\le \ConeToneVariable{x}\le \ConeToneVariable{b}.
\]

Let $\ConeToneVariable{d}=\ConeToneVariable{b}-\ConeToneVariable{a}$. Then
\[
\CfiveTonePopulationMean{\mu}-\ConeToneVariable{a}=\frac{\ConeToneVariable{2}\ConeToneVariable{d}}{3},
\qquad
\CfiveTonePopulationSD{\sigma}=\frac{\ConeToneVariable{d}}{\sqrt{18}},
\]
so
\[
\CfiveTonePopulationMean{\mu}-\CfiveTonePopulationSD{\sigma}-\ConeToneVariable{a}
=\ConeToneVariable{d}\left(\frac23-\frac{\ConeToneVariable{1}}{\sqrt{18}}\right),
\qquad
\CfiveTonePopulationMean{\mu}+\CfiveTonePopulationSD{\sigma}-\ConeToneVariable{a}
=\ConeToneVariable{d}\left(\frac23+\frac{\ConeToneVariable{1}}{\sqrt{18}}\right).
\]

Therefore,
\begin{align*}
	P(\mu-\sigma<X<\mu+\sigma)
	&=
	\int_{\mu-\sigma}^{\mu+\sigma}
	\frac{2(x-a)}{d^2}\,dx
	\\
	&=
	\left[
	\frac{(x-a)^2}{d^2}
	\right]_{\mu-\sigma}^{\mu+\sigma}
	\\
	&=
	\frac{(\mu+\sigma-a)^2-(\mu-\sigma-a)^2}{d^2}
	\\
	&=
	\left(\frac23+\frac{1}{\sqrt{18}}\right)^2
	-
	\left(\frac23-\frac{1}{\sqrt{18}}\right)^2
	\\
	&=
	\left[
	\left(\frac23\right)^2
	+2\left(\frac23\right)\frac{1}{\sqrt{18}}
	+\frac1{18}
	\right]
	-
	\left[
	\left(\frac23\right)^2
	-2\left(\frac23\right)\frac{1}{\sqrt{18}}
	+\frac1{18}
	\right]
	\\
	&=
	4\left(\frac23\right)\left(\frac{1}{\sqrt{18}}\right)
	\\
	&=
	\frac{8}{3\sqrt{18}}
	=
	\frac{8}{9\sqrt{2}}
	\approx 0.6285.
\end{align*}

For a linear decreasing distribution, the density is the reflection of the
increasing density about the midpoint of $[\ConeToneVariable{a},\ConeToneVariable{b}]$. Reflection preserves
probabilities and standardizes the distribution in the same way. Hence,
\[
\ConeToneProbability{P}(\CfiveTonePopulationMean{\mu}-\CfiveTonePopulationSD{\sigma}<\ConeToneVariable{X}<\CfiveTonePopulationMean{\mu}+\CfiveTonePopulationSD{\sigma})
=
\frac{\ConeToneVariable{8}}{3\sqrt{18}}
\approx \ConeToneProbability{62.85\%}.
\]

Thus, for every linear increasing or decreasing distribution on $[\ConeToneVariable{a},\ConeToneVariable{b}]$,
\[
\boxed{
	\ConeToneProbability{P}(\CfiveTonePopulationMean{\mu}-\CfiveTonePopulationSD{\sigma}<\ConeToneVariable{X}<\CfiveTonePopulationMean{\mu}+\CfiveTonePopulationSD{\sigma})
	=
	\frac{\ConeToneVariable{8}}{3\sqrt{18}}
	=
	\frac{\ConeToneVariable{8}}{9\sqrt{2}}
	\approx \ConeToneProbability{62.85\%}
}
\]

\subsection{Triangular distribution}\label{sec:c4-problem-triangular}

\SolutionStage{1}{Distribution and parameters}

A financial institution processes applications for small business loans.
Based on previous applications, the processing time ranges from 2 to
10 hours. Most applications can be completed in about 5 hours, while
shorter and longer processing times are less common. Therefore, the
processing time has a triangular probability density with a \CtwoToneMode{mode} at
5 hours.

Let $\ConeToneVariable{T}$ represent the processing time of a loan application, measured in
hours. Therefore,
\[
\ConeToneVariable{T}\sim\operatorname{Tri}(\ConeToneVariable{a},\ConeToneVariable{b},\CtwoToneMode{c})
\qquad \longrightarrow \qquad
\ConeToneVariable{T}\sim\operatorname{Tri}(2,10,5)
\qquad \longrightarrow \qquad
\ConeToneVariable{a}=\ConeToneVariable{2}\qquad \ConeToneVariable{b}=\ConeToneVariable{10}\qquad \CtwoToneMode{c}=\CtwoToneMode{5}
\]

\SolvedGraph{Triangular density}{
	\TriangularGraph{2}{5}{10}
}

\SolutionStage{2}{Mean}

For a triangular distribution,
\[
\CfiveToneExpected{E}(\ConeToneVariable{T})=\frac{\ConeToneVariable{a}+\ConeToneVariable{b}+\CtwoToneMode{c}}{3}
\qquad \longrightarrow \qquad
\CfiveToneExpected{E}(\ConeToneVariable{T})
=\frac{\ConeToneVariable{2}+\ConeToneVariable{10}+\ConeToneVariable{5}}{3}
=\frac{\ConeToneVariable{17}}{3}
\approx\CfiveToneExpected{5.67}
\qquad \longrightarrow \qquad
\boxed{\CfiveToneExpected{E}(\ConeToneVariable{T})\approx\CfiveToneExpected{5.67}\text{ hours}}
\]

\SolutionStage{3}{Variance}

For a triangular distribution,
\[
V\ConeToneVariable{a}\ConeToneVariable{r}(\ConeToneVariable{T})
=
\frac{a^2+b^2+c^2-\ConeToneVariable{a}\ConeToneVariable{b}-\ConeToneVariable{a}\CtwoToneMode{c}-\ConeToneVariable{b}\CtwoToneMode{c}}{18}.
\]

\[
\longrightarrow \qquad
\begin{aligned}
	V\ConeToneVariable{a}\ConeToneVariable{r}(\ConeToneVariable{T})
	&=
	\frac{
		\ConeToneVariable{2}^2+\ConeToneVariable{10}^2+\ConeToneVariable{5}^2
		-(\ConeToneVariable{2})(\ConeToneVariable{10})
		-(\ConeToneVariable{2})(\ConeToneVariable{5})
		-(\ConeToneVariable{10})(\ConeToneVariable{5})
	}{18}
	\\
	&=
	\frac{\ConeToneVariable{4}+\ConeToneVariable{100}+\ConeToneVariable{25}\ConeToneVariable{-20}\ConeToneVariable{-10}\ConeToneVariable{-50}}{18}
	\\
	&=
	\frac{\ConeToneVariable{49}}{18}
	\approx\ConeToneVariable{2.72}
\end{aligned}
\qquad \longrightarrow \qquad
\boxed{V\ConeToneVariable{a}\ConeToneVariable{r}(\ConeToneVariable{T})\approx\ConeToneVariable{2.72}\text{ hours}^2}
\]

\SolutionStage{4}{Interpretation}

The standard deviation is
\[
\CfiveTonePopulationSD{\sigma_T}=\sqrt{\frac{\ConeToneVariable{49}}{18}}
\approx\CfiveTonePopulationSD{1.65}
\qquad \longrightarrow \qquad
\boxed{\CfiveTonePopulationSD{\sigma_T}\approx\CfiveTonePopulationSD{1.65}\text{ hours}}.
\]

\SolutionStage{5}{Probability within one standard deviation}

We want the probability that the processing time is within one standard
deviation of the mean:
\[
\ConeToneProbability{P}(\CfiveTonePopulationMean{\mu}-\CfiveTonePopulationSD{\sigma_T}<\ConeToneVariable{T}<\CfiveTonePopulationMean{\mu}+\CfiveTonePopulationSD{\sigma_T})
\]

Using $\CfiveTonePopulationMean{\mu}\approx\CfiveTonePopulationMean{5.67}$ and $\CfiveTonePopulationSD{\sigma_T}\approx\CfiveTonePopulationSD{1.65}$,
\[
\ConeToneProbability{P}(5.67-1.65<\ConeToneVariable{T}<5.67+1.65)
=
\ConeToneProbability{P}(\ConeToneProbability{4.02}<\ConeToneVariable{T}<7.32)
\]

Since the \CtwoToneMode{mode} is $\CtwoToneMode{c}=\CtwoToneMode{5}$, the probability density is
\[
\CtwoToneDensity{f}(\ConeToneVariable{t})=
\begin{cases}
	\dfrac{\ConeToneVariable{2}(\ConeToneVariable{t}-\ConeToneVariable{a})}{(\ConeToneVariable{b}-\ConeToneVariable{a})(\CtwoToneMode{c}-\ConeToneVariable{a})}, & \ConeToneVariable{a}\leq \ConeToneVariable{t}\leq \CtwoToneMode{c},\\[\CtwoToneDensity{2}mm]
	\dfrac{\ConeToneVariable{2}(\ConeToneVariable{b}-\ConeToneVariable{t})}{(\ConeToneVariable{b}-\ConeToneVariable{a})(\ConeToneVariable{b}-\CtwoToneMode{c})}, & \CtwoToneMode{c}<\ConeToneVariable{t}\leq \ConeToneVariable{b}.
\end{cases}
\]

Therefore,
\[
\ConeToneProbability{P}(4.02<\ConeToneVariable{T}<7.32)
=
\int_{4.02}^{5}
\frac{\ConeToneVariable{2}(\ConeToneVariable{t}\ConeToneVariable{-2})}{(10-2)(5-2)}\,\ConeToneVariable{d}\ConeToneVariable{t}
+
\int_{5}^{7.32}
\frac{\ConeToneVariable{2}(\ConeToneVariable{10}-\ConeToneVariable{t})}{(10-2)(10-5)}\,\ConeToneVariable{d}\ConeToneVariable{t}
\]

\[
\approx0.563
\]

\[
\boxed{\ConeToneProbability{P}(\CfiveTonePopulationMean{\mu}-\CfiveTonePopulationSD{\sigma_T}<\ConeToneVariable{T}<\CfiveTonePopulationMean{\mu}+\CfiveTonePopulationSD{\sigma_T})\approx\ConeToneProbability{0.563}}
\]

Thus, approximately $56.3\%$ of loan-processing times are within one
standard deviation of the mean.


\subsection{Exponential distribution}\label{sec:c4-problem-exponential}

\SolutionStage{1}{Parameter}

An online clothing store receives orders throughout the day. Based on its
recent sales activity, the store receives an average of 15 orders per
hour. Let $\ConeToneVariable{X}$ represent the time between consecutive orders, measured in
minutes.

For an exponential distribution, the rate is the average number of events
per unit of time. Since 15 orders occur per hour,
\[
\CfourToneRate{\lambda}=\frac{\ConeToneVariable{15}}{60}
=\frac14\text{ min}^{-1}
\qquad \longrightarrow \qquad
\boxed{\CfourToneRate{\lambda}=\frac14\text{ min}^{-1}}
\]

\SolvedGraph{Exponential density}{
	\ExponentialGraph{0.25}
}

\SolutionStage{2}{Mean}

For an exponential distribution,
\[
\CfiveToneExpected{E}(\ConeToneVariable{X})=\frac{\ConeToneVariable{1}}{\CfourToneRate{\lambda}}
\qquad \longrightarrow \qquad
\CfiveToneExpected{E}(\ConeToneVariable{X})
=\frac{\ConeToneVariable{1}}{1/4}
=\CfiveToneExpected{4}
\qquad \longrightarrow \qquad
\boxed{\CfiveToneExpected{E}(\ConeToneVariable{X})=\CfiveToneExpected{4}\text{ minutes}}
\]

\SolutionStage{3}{Variance}

For an exponential distribution,
\[
V\ConeToneVariable{a}\ConeToneVariable{r}(\ConeToneVariable{X})=\frac{\ConeToneVariable{1}}{\CfourToneRate{\lambda^2}}
\qquad \longrightarrow \qquad
V\ConeToneVariable{a}\ConeToneVariable{r}(\ConeToneVariable{X})
=\frac{\ConeToneVariable{1}}{(1/4)^2}
=\ConeToneVariable{16}
\qquad \longrightarrow \qquad
\boxed{V\ConeToneVariable{a}\ConeToneVariable{r}(\ConeToneVariable{X})=\ConeToneVariable{16}\text{ minutes}^2}
\]

\SolutionStage{4}{Interpretation}

The standard deviation is
\[
\CfiveTonePopulationSD{\sigma_X}=\sqrt{\CfiveTonePopulationVariance{16}}=\CfiveTonePopulationSD{4}
\qquad \longrightarrow \qquad
\boxed{\CfiveTonePopulationSD{\sigma_X}=\CfiveTonePopulationSD{4}\text{ minutes}}
\]

The mean and standard deviation are both 4 minutes, reflecting the
characteristic variability of an exponential waiting-time distribution.

\SolutionStage{5}{Probability within one standard deviation}

We want the probability that the time between consecutive orders is
within one standard deviation of the mean:
\[
\ConeToneProbability{P}(\CfiveTonePopulationMean{\mu}-\CfiveTonePopulationSD{\sigma_X}<\ConeToneVariable{X}<\CfiveTonePopulationMean{\mu}+\CfiveTonePopulationSD{\sigma_X})
\]

Using $\CfiveTonePopulationMean{\mu}=\CfiveTonePopulationMean{4}$ and $\CfiveTonePopulationSD{\sigma_X}=\CfiveTonePopulationSD{4}$,
\[
\ConeToneProbability{P}(4-4<\ConeToneVariable{X}<4+4)
=
\ConeToneProbability{P}(\ConeToneProbability{0}<\ConeToneVariable{X}<8)
\]

For an exponential distribution, the \ConeToneCDF{cumulative} \ConeToneCDF{distribution} \ConeToneCDF{function} is
\[
\ConeToneProbability{P}(\ConeToneVariable{X}<\ConeToneVariable{x})=\ConeToneProbability{1}-e^{-\lambda x}.
\]

Therefore,
\[
\ConeToneProbability{P}(0<\ConeToneVariable{X}<8)
=
\ConeToneProbability{1}-e^{-(1/4)(8)}
=
\ConeToneProbability{1}-e^{-2}
\approx\ConeToneProbability{0.8647}
\]

\[
\boxed{\ConeToneProbability{P}(\CfiveTonePopulationMean{\mu}-\CfiveTonePopulationSD{\sigma_X}<\ConeToneVariable{X}<\CfiveTonePopulationMean{\mu}+\CfiveTonePopulationSD{\sigma_X})\approx\ConeToneProbability{0.865}}
\]

Thus, approximately $86.5\%$ of the time intervals between consecutive
orders are within one standard deviation of the mean.


\subsection{Normal distribution}\label{sec:c4-problem-normal}

\SolutionStage{1}{Parameters}

A small retail business tracks its monthly sales revenue over several
months. Based on its historical sales data, the monthly revenue is
approximately centered at \$45,000, with a standard deviation of
\$8,000. Let $\ConeToneVariable{Y}$ represent the monthly sales revenue, measured in
thousands of dollars.

Therefore,
\[
\ConeToneVariable{Y}\sim N(45,8^2)
\qquad \longrightarrow \qquad
\CfiveTonePopulationMean{\mu}=\CfiveTonePopulationMean{45}\qquad \CfiveTonePopulationSD{\sigma}=\CfiveTonePopulationSD{8}
\]

\SolvedGraph{Normal density}{
	\NormalGraph{45}{8}
}

\SolutionStage{2}{Mean}

For a normal distribution,
\[
\CfiveToneExpected{E}(\ConeToneVariable{Y})=\CfiveTonePopulationMean{\mu}
\qquad \longrightarrow \qquad
\boxed{\CfiveToneExpected{E}(\ConeToneVariable{Y})=\CfiveToneExpected{45}\text{ thousand dollars}}
\]

\SolutionStage{3}{Variance}

For a normal distribution,
\[
V\ConeToneVariable{a}\ConeToneVariable{r}(\ConeToneVariable{Y})=\CfiveTonePopulationVariance{\sigma^2}
\qquad \longrightarrow \qquad
V\ConeToneVariable{a}\ConeToneVariable{r}(\ConeToneVariable{Y})=\ConeToneVariable{8}^2=\ConeToneVariable{64}
\qquad \longrightarrow \qquad
\boxed{V\ConeToneVariable{a}\ConeToneVariable{r}(\ConeToneVariable{Y})=\ConeToneVariable{64}\text{ thousand dollars}^2}.
\]

\SolutionStage{4}{Interpretation}

The standard deviation is
\[
\CfiveTonePopulationSD{\sigma_Y}=\sqrt{\CfiveTonePopulationVariance{64}}=\CfiveTonePopulationSD{8}
\qquad \longrightarrow \qquad
\boxed{\CfiveTonePopulationSD{\sigma_Y}=\CfiveTonePopulationSD{8}\text{ thousand dollars}}.
\]

This means that monthly sales revenue typically varies by about
\$8,000 around the mean revenue of \$45,000.

\SolutionStage{5}{Probability within one standard deviation}

We want the probability that monthly sales revenue is within one standard
deviation of the mean:
\[
\ConeToneProbability{P}(\CfiveTonePopulationMean{\mu}-\CfiveTonePopulationSD{\sigma_Y}<\ConeToneVariable{Y}<\CfiveTonePopulationMean{\mu}+\CfiveTonePopulationSD{\sigma_Y})
\]

Using $\CfiveTonePopulationMean{\mu}=\CfiveTonePopulationMean{45}$ and $\CfiveTonePopulationSD{\sigma_Y}=\CfiveTonePopulationSD{8}$,
\[
\ConeToneProbability{P}(45-8<\ConeToneVariable{Y}<45+8)
=
\ConeToneProbability{P}(\ConeToneProbability{37}<\ConeToneVariable{Y}<53)
\]

Standardizing the two endpoints,
\[
\CeightToneCritical{z_1}=\frac{\ConeToneVariable{37}\ConeToneVariable{-45}}{8}=\CeightToneCritical{-1},
\qquad
\CeightToneCritical{z_2}=\frac{\ConeToneVariable{53}\ConeToneVariable{-45}}{8}=\CeightToneCritical{1}
\]

Therefore,
\[
\ConeToneProbability{P}(37<\ConeToneVariable{Y}<53)
=
\ConeToneProbability{P}(\ConeToneProbability{-1}<\ConeToneVariable{Z}<1)
\approx\ConeToneProbability{0.6827}
\]

\[
\boxed{\ConeToneProbability{P}(\CfiveTonePopulationMean{\mu}-\CfiveTonePopulationSD{\sigma_Y}<\ConeToneVariable{Y}<\CfiveTonePopulationMean{\mu}+\CfiveTonePopulationSD{\sigma_Y})\approx\ConeToneProbability{0.683}}
\]

Thus, approximately $68.3\%$ of monthly sales revenues are within one
standard deviation of the mean.


\newpage

\Needspace{8\baselineskip}
\subsection*{Notation reference}
\phantomsection\label{tone:C5:reference}
\COneToneDetailedReference{C5}
\section{Context Problems}\label{sec:c4-context-problems}

In each problem, the distribution is not named. Use the contextual clues to
identify the underlying distribution and its parameters. Then calculate the
mean, variance, and standard deviation, and interpret the mean and standard
deviation in context.

\SectionProblemsTableOfContents{
	\SectionProblemLink{sec:c4-context-shuttle}{Airport shuttle arrival} \\
	\SectionProblemLink{sec:c4-context-irrigation}{Irrigation duration} \\
	\SectionProblemLink{sec:c4-context-craft}{Craft project time} \\
	\SectionProblemLink{sec:c4-context-auction}{Online auction price} \\
	\SectionProblemLink{sec:c4-context-support}{Technical support time} \\
	\SectionProblemLink{sec:c4-context-battery}{Battery replacement age} \\
	\SectionProblemLink{sec:c4-context-repair}{Bicycle repair time} \\
	\SectionProblemLink{sec:c4-context-loading}{Truck loading time} \\
	\SectionProblemLink{sec:c4-context-checkout}{Checkout arrivals} \\
	\SectionProblemLink{sec:c4-context-breakdown}{Machine breakdowns} \\
	\SectionProblemLink{sec:c4-context-temperature}{Greenhouse temperature} \\
	\SectionProblemLink{sec:c4-context-filling}{Package filling mass}
}

\subsection{Airport shuttle arrival}\label{sec:c4-context-shuttle}
An airport shuttle may arrive at any time from 8 to 20 minutes after a
passenger reaches the stop, and every time in that interval is equally likely.
Let $\ConeToneVariable{X}$ be the waiting time in minutes.

\textbf{Task:} Identify the distribution and its parameters, then find
$\CfiveToneExpected{E}(\ConeToneVariable{X})$, $V\ConeToneVariable{a}\ConeToneVariable{r}(\ConeToneVariable{X})$, and $\CfiveTonePopulationSD{\sigma_X}$.

\SolutionStage{1}{Solution}
The phrase ``every time in that interval is equally likely'' identifies
$\ConeToneVariable{X}\sim \ConeToneVariable{U}(8,20)$. Therefore,
\[
\CfiveToneExpected{E}(\ConeToneVariable{X})=\frac{\ConeToneVariable{8}+\ConeToneVariable{20}}{2}=\CfiveToneExpected{14},\qquad
V\ConeToneVariable{a}\ConeToneVariable{r}(\ConeToneVariable{X})=\frac{(\ConeToneVariable{20}\ConeToneVariable{-8})^2}{12}=\ConeToneVariable{12},\qquad
\CfiveTonePopulationSD{\sigma_X}=\sqrt{\CfiveTonePopulationVariance{12}}\approx\CfiveTonePopulationSD{3.46}.
\]
The average wait is $14$ minutes, with a typical distance from that average of
about $3.46$ minutes.

\subsection{Irrigation duration}\label{sec:c4-context-irrigation}
A field's automatic irrigation cycle lasts between 30 and 54 minutes. A sensor
selects the stopping time so that every duration in this interval is equally
likely. Let $\ConeToneVariable{X}$ be the cycle duration in minutes.

\textbf{Task:} Identify the distribution and its parameters, then find
$\CfiveToneExpected{E}(\ConeToneVariable{X})$, $V\ConeToneVariable{a}\ConeToneVariable{r}(\ConeToneVariable{X})$, and $\CfiveTonePopulationSD{\sigma_X}$.

\SolutionStage{1}{Solution}
Equal likelihood throughout a fixed interval identifies $\ConeToneVariable{X}\sim \ConeToneVariable{U}(30,54)$.
Thus,
\[
\CfiveToneExpected{E}(\ConeToneVariable{X})=\CfiveToneExpected{42},\qquad V\ConeToneVariable{a}\ConeToneVariable{r}(\ConeToneVariable{X})=\frac{\ConeToneVariable{24}^2}{12}=\ConeToneVariable{48},\qquad
\CfiveTonePopulationSD{\sigma_X}=\sqrt{\CfiveTonePopulationVariance{48}}\approx\CfiveTonePopulationSD{6.93}.
\]
Cycles average $42$ minutes and typically differ from that average by about
$6.93$ minutes.

\subsection{Craft project time}\label{sec:c4-context-craft}
A beginner needs between 2 and 8 hours to finish a craft project. Longer times
are progressively more likely, and the probability density rises in a straight
line from zero at 2 hours to its maximum at 8 hours. Let $\ConeToneVariable{X}$ be the completion
time in hours.

\textbf{Task:} Identify the distribution and its parameters, then find
$\CfiveToneExpected{E}(\ConeToneVariable{X})$, $V\ConeToneVariable{a}\ConeToneVariable{r}(\ConeToneVariable{X})$, and $\CfiveTonePopulationSD{\sigma_X}$.

\SolutionStage{1}{Solution}
A density that rises linearly from the lower endpoint identifies
$\ConeToneVariable{X}\sim\operatorname{Lin}_{\mathrm{inc}}(2,8)$. Hence,
\[
\CfiveToneExpected{E}(\ConeToneVariable{X})=\frac{\ConeToneVariable{2}+\ConeToneVariable{2}(\ConeToneVariable{8})}{3}=\CfiveToneExpected{6},\qquad
V\ConeToneVariable{a}\ConeToneVariable{r}(\ConeToneVariable{X})=\frac{(\ConeToneVariable{8}\ConeToneVariable{-2})^2}{18}=\ConeToneVariable{2},\qquad
\CfiveTonePopulationSD{\sigma_X}=\sqrt\CfiveTonePopulationSD{2}\approx\CfiveTonePopulationSD{1.41}.
\]
The mean completion time is $6$ hours, with a standard deviation of about
$1.41$ hours.

\subsection{Online auction price}\label{sec:c4-context-auction}
The final price of a collectible lies between \$10 and \$25. Within this range,
higher prices become progressively more likely, and the density increases in a
straight line from zero at \$10 to its maximum at \$25. Let $\ConeToneVariable{X}$ be the price in
dollars.

\textbf{Task:} Identify the distribution and its parameters, then find
$\CfiveToneExpected{E}(\ConeToneVariable{X})$, $V\ConeToneVariable{a}\ConeToneVariable{r}(\ConeToneVariable{X})$, and $\CfiveTonePopulationSD{\sigma_X}$.

\SolutionStage{1}{Solution}
The contextual shape identifies
$\ConeToneVariable{X}\sim\operatorname{Lin}_{\mathrm{inc}}(10,25)$. Therefore,
\[
\CfiveToneExpected{E}(\ConeToneVariable{X})=\frac{\ConeToneVariable{10}+\ConeToneVariable{2}(\ConeToneVariable{25})}{3}=\CfiveToneExpected{20},\qquad
V\ConeToneVariable{a}\ConeToneVariable{r}(\ConeToneVariable{X})=\frac{\ConeToneVariable{15}^2}{18}=\ConeToneVariable{12.5},\qquad
\CfiveTonePopulationSD{\sigma_X}=\sqrt{\CfiveTonePopulationVariance{12.5}}\approx\CfiveTonePopulationSD{3.54}.
\]
The mean final price is \$20, with a standard deviation of about \$3.54.

\subsection{Technical support time}\label{sec:c4-context-support}
A simple technical-\ConeToneVariable{support} call lasts between 0 and 30 minutes. Very short calls
are most likely, and the probability density decreases in a straight line to
zero at 30 minutes. Let $\ConeToneVariable{X}$ be the call duration in minutes.

\textbf{Task:} Identify the distribution and its parameters, then find
$\CfiveToneExpected{E}(\ConeToneVariable{X})$, $V\ConeToneVariable{a}\ConeToneVariable{r}(\ConeToneVariable{X})$, and $\CfiveTonePopulationSD{\sigma_X}$.

\SolutionStage{1}{Solution}
A density that falls linearly to the upper endpoint identifies
$\ConeToneVariable{X}\sim\operatorname{Lin}_{\mathrm{dec}}(0,30)$. Thus,
\[
\CfiveToneExpected{E}(\ConeToneVariable{X})=\frac{\ConeToneVariable{2}(\ConeToneVariable{0})+\ConeToneVariable{30}}{3}=\CfiveToneExpected{10},\qquad
V\ConeToneVariable{a}\ConeToneVariable{r}(\ConeToneVariable{X})=\frac{\ConeToneVariable{30}^2}{18}=\ConeToneVariable{50},\qquad
\CfiveTonePopulationSD{\sigma_X}=\sqrt{\CfiveTonePopulationVariance{50}}\approx\CfiveTonePopulationSD{7.07}.
\]
Calls average $10$ minutes and have a standard deviation of about $\CfiveTonePopulationSD{7.07}$
minutes.

\subsection{Battery replacement age}\label{sec:c4-context-battery}
A fleet replaces a certain battery between 12 and 48 months after installation.
Replacement is most likely near 12 months, and the density decreases linearly
to zero at 48 months. Let $\ConeToneVariable{X}$ be the replacement age in months.

\textbf{Task:} Identify the distribution and its parameters, then find
$\CfiveToneExpected{E}(\ConeToneVariable{X})$, $V\ConeToneVariable{a}\ConeToneVariable{r}(\ConeToneVariable{X})$, and $\CfiveTonePopulationSD{\sigma_X}$.

\SolutionStage{1}{Solution}
These clues identify $\ConeToneVariable{X}\sim\operatorname{Lin}_{\mathrm{dec}}(12,48)$. Hence,
\[
\CfiveToneExpected{E}(\ConeToneVariable{X})=\frac{\ConeToneVariable{2}(\ConeToneVariable{12})+\ConeToneVariable{48}}{3}=\CfiveToneExpected{24},\qquad
V\ConeToneVariable{a}\ConeToneVariable{r}(\ConeToneVariable{X})=\frac{\ConeToneVariable{36}^2}{18}=\ConeToneVariable{72},\qquad
\CfiveTonePopulationSD{\sigma_X}=\sqrt{\CfiveTonePopulationVariance{72}}\approx\CfiveTonePopulationSD{8.49}.
\]
The mean replacement age is $24$ months, with a standard deviation of about
$8.49$ months.

\subsection{Bicycle repair time}\label{sec:c4-context-repair}
A bicycle repair takes between 20 and 50 minutes. Repairs near 32 minutes are
most common, and times become linearly less likely on either side of that peak.
Let $\ConeToneVariable{X}$ be the repair time in minutes.

\textbf{Task:} Identify the distribution and its parameters, then find
$\CfiveToneExpected{E}(\ConeToneVariable{X})$, $V\ConeToneVariable{a}\ConeToneVariable{r}(\ConeToneVariable{X})$, and $\CfiveTonePopulationSD{\sigma_X}$.

\SolutionStage{1}{Solution}
A bounded density with one peak and straight sides identifies
$\ConeToneVariable{X}\sim\operatorname{Tri}(20,50,32)$. Therefore,
\[
\CfiveToneExpected{E}(\ConeToneVariable{X})=\frac{\ConeToneVariable{20}+\ConeToneVariable{50}+\ConeToneVariable{32}}{3}=\CfiveToneExpected{34},
\]
\[
V\ConeToneVariable{a}\ConeToneVariable{r}(\ConeToneVariable{X})=\frac{\ConeToneVariable{20}^2+\ConeToneVariable{50}^2+\ConeToneVariable{32}^2-(\ConeToneVariable{20})(\ConeToneVariable{50})-(\ConeToneVariable{20})(\ConeToneVariable{32})-(\ConeToneVariable{50})(\ConeToneVariable{32})}{18}=\ConeToneVariable{38},
\qquad \CfiveTonePopulationSD{\sigma_X}=\sqrt{\CfiveTonePopulationVariance{38}}\approx\CfiveTonePopulationSD{6.16}.
\]
Repairs average $34$ minutes and typically differ from that average by about
$6.16$ minutes.

\subsection{Truck loading time}\label{sec:c4-context-loading}
Loading a delivery truck takes between 5 and 17 minutes. The most likely time is
8 minutes, and the density changes linearly from each endpoint to that peak.
Let $\ConeToneVariable{X}$ be the loading time in minutes.

\textbf{Task:} Identify the distribution and its parameters, then find
$\CfiveToneExpected{E}(\ConeToneVariable{X})$, $V\ConeToneVariable{a}\ConeToneVariable{r}(\ConeToneVariable{X})$, and $\CfiveTonePopulationSD{\sigma_X}$.

\SolutionStage{1}{Solution}
The minimum, maximum, and single linearly connected peak identify
$\ConeToneVariable{X}\sim\operatorname{Tri}(5,17,8)$. Thus,
\[
\CfiveToneExpected{E}(\ConeToneVariable{X})=\frac{\ConeToneVariable{5}+\ConeToneVariable{17}+\ConeToneVariable{8}}{3}=\CfiveToneExpected{10},
\]
\[
V\ConeToneVariable{a}\ConeToneVariable{r}(\ConeToneVariable{X})=\frac{\ConeToneVariable{5}^2+\ConeToneVariable{17}^2+\ConeToneVariable{8}^2-(\ConeToneVariable{5})(\ConeToneVariable{17})-(\ConeToneVariable{5})(\ConeToneVariable{8})-(\ConeToneVariable{17})(\ConeToneVariable{8})}{18}=\ConeToneVariable{6.5},
\qquad \CfiveTonePopulationSD{\sigma_X}=\sqrt{\CfiveTonePopulationVariance{6.5}}\approx\CfiveTonePopulationSD{2.55}.
\]
The mean loading time is $10$ minutes, with a standard deviation of about
$2.55$ minutes.

\subsection{Checkout arrivals}\label{sec:c4-context-checkout}
Customers arrive independently at a checkout at a steady average rate of six
customers per hour. Let $\ConeToneVariable{X}$ be the time in minutes between consecutive
arrivals.

\textbf{Task:} Identify the distribution and its parameters, then find
$\CfiveToneExpected{E}(\ConeToneVariable{X})$, $V\ConeToneVariable{a}\ConeToneVariable{r}(\ConeToneVariable{X})$, and $\CfiveTonePopulationSD{\sigma_X}$.

\SolutionStage{1}{Solution}
A waiting time between independent events occurring at a steady rate identifies
an exponential model. In minutes, $\CfourToneRate{\lambda}=\CfourToneRate{6}/60=\CfourToneRate{0.1}$, so
$\ConeToneVariable{X}\sim\operatorname{Exp}(0.1)$. Therefore,
\[
\CfiveToneExpected{E}(\ConeToneVariable{X})=\frac1{0.1}=\CfiveToneExpected{10},\qquad
V\ConeToneVariable{a}\ConeToneVariable{r}(\ConeToneVariable{X})=\frac1{0.1^2}=\ConeToneVariable{100},\qquad
\CfiveTonePopulationSD{\sigma_X}=\CfiveTonePopulationSD{10}.
\]
The mean and standard deviation of the time between arrivals are both
$10$ minutes.

\subsection{Machine breakdowns}\label{sec:c4-context-breakdown}
Breakdowns of a machine occur independently at a steady average rate of four
per hour. Let $\ConeToneVariable{X}$ be the time in hours from one breakdown to the next.

\textbf{Task:} Identify the distribution and its parameters, then find
$\CfiveToneExpected{E}(\ConeToneVariable{X})$, $V\ConeToneVariable{a}\ConeToneVariable{r}(\ConeToneVariable{X})$, and $\CfiveTonePopulationSD{\sigma_X}$.

\SolutionStage{1}{Solution}
The waiting-time clues identify $\ConeToneVariable{X}\sim\operatorname{Exp}(4)$. Hence,
\[
\CfiveToneExpected{E}(\ConeToneVariable{X})=\frac14=\CfiveToneExpected{0.25},\qquad
V\ConeToneVariable{a}\ConeToneVariable{r}(\ConeToneVariable{X})=\frac1{4^2}=\ConeToneVariable{0.0625},\qquad
\CfiveTonePopulationSD{\sigma_X}=\CfiveTonePopulationSD{0.25}.
\]
The mean time between breakdowns is $0.25$ hour (15 minutes), and its standard
deviation is also $0.25$ hour.

\subsection{Greenhouse temperature}\label{sec:c4-context-temperature}
Daily noon temperatures in a controlled greenhouse form an approximately
symmetric, bell-shaped pattern centered at $68^\circ$F, with a standard
deviation of $7^\circ$F. Let $\ConeToneVariable{X}$ be the temperature in degrees Fahrenheit.

\textbf{Task:} Identify the distribution and its parameters, then find
$\CfiveToneExpected{E}(\ConeToneVariable{X})$, $V\ConeToneVariable{a}\ConeToneVariable{r}(\ConeToneVariable{X})$, and $\CfiveTonePopulationSD{\sigma_X}$.

\SolutionStage{1}{Solution}
The symmetric bell shape identifies $\ConeToneVariable{X}\sim N(68,7^2)$. Thus,
\[
\CfiveToneExpected{E}(\ConeToneVariable{X})=\CfiveToneExpected{68},\qquad V\ConeToneVariable{a}\ConeToneVariable{r}(\ConeToneVariable{X})=\ConeToneVariable{7}^2=\ConeToneVariable{49},\qquad \CfiveTonePopulationSD{\sigma_X}=\CfiveTonePopulationSD{7}.
\]
The mean temperature is $68^\circ$F, and temperatures typically differ from it
by about $7^\circ$F.

\subsection{Package filling mass}\label{sec:c4-context-filling}
A calibrated machine fills packages with masses that follow an approximately
symmetric, bell-shaped pattern centered at 500 grams. The standard deviation is
25 grams. Let $\ConeToneVariable{X}$ be the mass of a package in grams.

\textbf{Task:} Identify the distribution and its parameters, then find
$\CfiveToneExpected{E}(\ConeToneVariable{X})$, $V\ConeToneVariable{a}\ConeToneVariable{r}(\ConeToneVariable{X})$, and $\CfiveTonePopulationSD{\sigma_X}$.

\SolutionStage{1}{Solution}
The contextual shape identifies $\ConeToneVariable{X}\sim N(500,25^2)$. Therefore,
\[
\CfiveToneExpected{E}(\ConeToneVariable{X})=\CfiveToneExpected{500},\qquad V\ConeToneVariable{a}\ConeToneVariable{r}(\ConeToneVariable{X})=\ConeToneVariable{25}^2=\ConeToneVariable{625},\qquad \CfiveTonePopulationSD{\sigma_X}=\CfiveTonePopulationSD{25}.
\]
The mean filling mass is $500$ grams, with a standard deviation of $\CfiveTonePopulationSD{25}$ grams.


\end{document}


